\section*{Problem \#306}

\subsection*{1) ROUND-2 OBJECTIVE}
\textbf{Path C (obstruction/correction).}
The full statement remains open, so the realistic round-2 target is a corrected proven statement covering genuinely new infinite subclasses.
In this round I prove two such subclasses:
\begin{enumerate}
\item a graph-theoretic construction for many fractions whose denominator is a product of distinct primes;
\item a short explicit three-term criterion for certain unit fractions \(1/p\) with \(p\) prime.
\end{enumerate}
These are strictly stronger than the isolated examples from Round~1.

\subsection*{2) Round-1 FOUNDATION USED}
I use the following Round~1 facts.
\begin{enumerate}
\item If a rational number is representable as a sum of reciprocals of distinct semiprimes, then its reduced denominator is squarefree.
\item Round~1 recorded explicit semiprime decompositions for \(1/3,1/5,1/2,2/5\), and a 48-term semiprime decomposition of \(1\).
\end{enumerate}

\subsection*{3) NEW INSIGHT / TOOL (ROUND-2)}
The new insight is that the allowed denominators \(p_i p_j\) naturally correspond to edges of a graph on the prime factors of the target denominator.
This gives an immediate exact formula for a large family of fractions.
A second new observation is that for prime \(p\), a three-term semiprime decomposition of \(1/p\) appears whenever
\[
  (q-1)(r-1)=p+1
\]
for suitable distinct primes \(q,r\).

\subsection*{4) ATTACK PLAN (ROUND-2)}
The specific gap after Round~1 was the lack of any systematic constructive theorem beyond a few hand-picked examples.
To close that gap I will:
\begin{enumerate}
\item prove a general graph formula for denominators \(P=\prod p_i\);
\item extract from it a concrete infinite family for three-prime denominators;
\item prove a separate three-term criterion for \(1/p\) and show that it subsumes several Round~1 examples.
\end{enumerate}

\subsection*{5) WORK (ROUND-2)}
\textbf{Theorem A (graph construction).}
Let \(p_1,\dots,p_t\) be distinct primes, and let
\[
  P:=p_1p_2\cdots p_t.
\]
For any simple graph \(G\) on the vertex set \(\{1,\dots,t\}\), define
\[
  a_G:=\sum_{\{i,j\}\in E(G)} \prod_{k\ne i,j} p_k.
\]
Then
\[
  \frac{a_G}{P}=\sum_{\{i,j\}\in E(G)} \frac{1}{p_i p_j}.
\]
In particular, every such \(a_G/P\) has a semiprime Egyptian-fraction decomposition with distinct denominators.

\medskip
\textit{Proof.}
For each edge \(\{i,j\}\),
\[
  \frac{1}{p_i p_j}=\frac{1}{P}\prod_{k\ne i,j} p_k.
\]
Summing over all edges of \(G\) gives exactly the claimed identity.
Since different edges give different unordered pairs \(\{i,j\}\), the denominators \(p_i p_j\) are distinct semiprimes.
\hfill \(\square\)

\medskip
\textbf{Corollary A1 (three-prime denominator family).}
Let \(p,q,r\) be distinct primes.
Then each of the seven fractions
\[
  \frac{p}{pqr},\ \frac{q}{pqr},\ \frac{r}{pqr},\ \frac{p+q}{pqr},\ \frac{p+r}{pqr},\ \frac{q+r}{pqr},\ \frac{p+q+r}{pqr}
\]
has a semiprime Egyptian-fraction decomposition.
Indeed,
\[
  \frac{p}{pqr}=\frac1{qr},
  \qquad
  \frac{q}{pqr}=\frac1{pr},
  \qquad
  \frac{r}{pqr}=\frac1{pq},
\]
\[
  \frac{p+q}{pqr}=\frac1{pr}+\frac1{qr},
  \qquad
  \frac{p+r}{pqr}=\frac1{pq}+\frac1{qr},
  \qquad
  \frac{q+r}{pqr}=\frac1{pq}+\frac1{pr},
\]
\[
  \frac{p+q+r}{pqr}=\frac1{pq}+\frac1{pr}+\frac1{qr}.
\]
This gives an explicit infinite family of squarefree denominators with three prime factors for which several nontrivial numerators are provably representable.

\medskip
\textbf{Theorem B (a three-term criterion for \(1/p\)).}
Let \(p,q,r\) be distinct primes.
If
\[
  (q-1)(r-1)=p+1,
\]
then
\[
  \frac1p=\frac1{pq}+\frac1{pr}+\frac1{qr}.
\]
Equivalently, whenever
\[
  qr-q-r=p,
\]
we obtain a three-term semiprime Egyptian-fraction decomposition of \(1/p\).

\medskip
\textit{Proof.}
The condition \((q-1)(r-1)=p+1\) expands to
\[
  qr-q-r=p.
\]
Therefore
\[
  \frac1{pq}+\frac1{pr}+\frac1{qr}
  =\frac{p+q+r}{pqr}
  =\frac{qr}{pqr}
  =\frac1p.
\]
The denominators are distinct semiprimes because \(p,q,r\) are distinct primes.
\hfill \(\square\)

\medskip
\textbf{Examples from Theorem B.}
\begin{enumerate}
\item If \(p+2\) is prime, then taking \(q=2\) and \(r=p+2\) gives
\[
  \frac1p=\frac1{2p}+\frac1{p(p+2)}+\frac1{2(p+2)}.
\]
Thus
\[
  \frac1{17}=\frac1{34}+\frac1{38}+\frac1{323}.
\]
\item If \((p+3)/2\) is prime, then taking \(q=3\) and \(r=(p+3)/2\) gives
\[
  \frac1p=\frac1{3p}+\frac1{p(p+3)/2}+\frac1{3(p+3)/2}.
\]
Thus
\[
  \frac1{11}=\frac1{21}+\frac1{33}+\frac1{77},
  \qquad
  \frac1{19}=\frac1{33}+\frac1{57}+\frac1{209}.
\]
\item Taking \((p,q,r)=(23,5,7)\) gives
\[
  \frac1{23}=\frac1{35}+\frac1{115}+\frac1{161}.
\]
\end{enumerate}
These formulas subsume several of the isolated small examples from Round~1 and show that there is a genuine infinite-pattern mechanism behind them.

\subsection*{6) ADVERSARIAL VERIFICATION}
\begin{enumerate}
\item \textbf{Are the denominators really distinct in Theorem A?}
Yes. Different edges correspond to different unordered pairs \(\{i,j\}\), and for distinct primes the products \(p_i p_j\) are distinct.
\item \textbf{Is Theorem A compatible with the original squarefree-denominator hypothesis?}
Yes. The target denominator before reduction is \(P=\prod p_i\), which is squarefree; after reduction, the denominator still divides \(P\), so it remains squarefree.
\item \textbf{Could Theorem B fail because one denominator repeats another?}
No. The theorem explicitly assumes \(p,q,r\) are distinct primes, so \(pq,pr,qr\) are pairwise distinct semiprimes.
\item \textbf{Did I accidentally claim necessity in Theorem B?}
No. The theorem is only a sufficient criterion. It does \emph{not} classify all three-term semiprime decompositions of \(1/p\).
\item \textbf{Does any of this solve the full problem?}
No. The full question asks for \emph{every} rational with squarefree denominator. The new theorems produce large explicit subclasses, but not all of them.
\end{enumerate}

\subsection*{7) FINAL (EXACTLY ONE)}
\textbf{UNRESOLVED (BUT STRICTLY ADVANCED).}
Round~2 proves two new constructive theorems.
First, the graph construction gives a systematic infinite family of representable rationals with squarefree denominator \(P=\prod p_i\).
Second, Theorem~B gives a clean explicit three-term criterion for many prime-denominator unit fractions \(1/p\).
This is a genuine step beyond the Round~1 list of isolated examples, but it still falls short of the full conjecture for arbitrary \(a/b\).

\subsection*{8) COMPLETION ESTIMATE}
\textbf{COMPLETION: 34\%}

\subsection*{9) REFERENCES}
\begin{itemize}
\item Round~1 output in this conversation (used as fixed foundation).
\item T. F. Bloom, \emph{Erd\H{o}s Problem \#306}.
\item S. Butler, P. Erd\H{o}s and R. Graham, \emph{Egyptian fractions with each denominator having three distinct prime divisors}.
\item T. Watanabe, \emph{New examples of the representation of \(1\) by the sum of reciprocals of semiprime numbers}, arXiv:2009.03275.
\end{itemize}
